\section{Three ontological readings (R1, R2, R3)}
\label{sec:ontological_readings_en}

PT, by its nature as a purely informational theory derived from
a single arithmetic parameter $s = 1/2$, admits three ontological
readings compatible with its formal structure. The trichotomy is
laid out in detail in chapters~24 and~26 of the
monograph~\cite{PT_Monograph}; we recall it here briefly and
apply it specifically to the Aharonov--Bohm / Bell debate
developed in this article. The crucial point to emphasise --- and
the substance of~\S~\ref{ssec:reading_independence_en} --- is
that the proofs of
Sections~\ref{sec:field_connection}--\ref{sec:qh1a_prediction}
are logically independent of the choice of reading. The
trichotomy bears on the \emph{interpretation} of what the
theorems say, not on their validity.

\subsection{R1: PT as mathematical coincidence}
\label{ssec:R1_en}

Reading R1 holds that PT reproduces the correct empirical content
--- the Aharonov--Bohm phase, the EPR correlations, the absence
of geometric dependence of the MI --- by a structural coincidence
between the combinatorics of the Eratosthenes sieve and the
phenomenological structure of quantum electrodynamics. No
ontological commitment is made on the physical nature of the
angles $\theta_p$: they are mathematical quantities useful for
prediction, and nothing more. The theorems of the preceding
sections retain their validity --- they are mathematical
assertions about the structure of the sieve --- but their
correspondence with observed phenomena is treated as an
unexplained empirical fact.

This reading is the epistemologically most cautious: it commits
only the minimum required to account for the observations. Its
weakness is that it offers no \emph{explanation} of the
coincidence: why does the combinatorics of the sieve yield
exactly the right phenomenological structure? R1 does not answer
this question.

\subsection{R2: structuralism}
\label{ssec:R2_en}

Reading R2 holds that the Eratosthenes sieve and quantum
electrodynamics share an abstract holonomic \emph{pattern}, and
that neither is ontologically prior. The $\theta_p$ are then
neither mathematical objects (in the sense of R1) nor physical
objects (in the sense of R3): they are components of an abstract
structure of which the two theories are incarnations.

This reading, classical in philosophy of science under the name
\emph{structuralism}~\cite{Belot1998,Maudlin2011}, finds in PT a
particularly favourable ground: the correspondence between the
sieve and the phenomenology is no longer a coincidence but the
expression of a common underlying structure, and the primacy of
holonomy becomes a structural property rather than a matter of
material ontology. The theorems of the preceding sections retain
their validity and acquire explanatory status in the sense that
they describe the common structure.

\subsection{R3: informational realism}
\label{ssec:R3_en}

Reading R3 holds that the $\theta_p$ \emph{are} the fundamental
physical reality, and that the classical objects of
electrodynamics --- $A_\mu$, $F_{\mu\nu}$, $g_{\mu\nu}$ --- are
phenomenological projections onto the emergent spacetime. Under
this reading, the ontological substance of the theory is
information ($\DKL$ in bits); the notions of particle, field,
space, and time are emergent constructions above this
informational substance.

R3 is the most committed reading: it asserts that the world is
made of arithmetic information, that matter and space are mere
projections of it, and that the correlations observed in Bell
tests are direct reflections of the fundamental CRT structure.
Under R3, the dissolution of spatial non-locality
(\S~\ref{ssec:dissolution_consequence_en}) takes its strongest
sense: the EPR correlations are literally non-spatial, because
what is measured has no spatial support.

R3 inherits the classical philosophical difficulties of
informational realism~\cite{Healey2007}: the nature of the
``execution'' of the sieve remains open (who or what
``computes'' the sequence $\{g_n\}$?), and the apparent
correspondence between the mathematical structure and the
observed phenomenology demands, under R3, a metaphysical
explanation that exceeds the scope of the theory itself.

\subsection{Theorem independence from the reading}
\label{ssec:reading_independence_en}

The central logical point of this section is that
Theorems~\ref{thm:invariance_geom_en}
and~\ref{thm:ruelle_zero_en},
Corollary~\ref{cor:closed_form_en}, and
Prediction~\ref{pred:qh1a_en} are \emph{all} valid under R1, R2,
R3. No formal result of the preceding sections depends on the
choice of ontological interpretation. The proofs are mathematical;
the predictions are functions of holonomy angles and the
emergent metric; and the numerical values (e.g.\
$\mathcal{I}_{3,5} = 0.138$ bits) are computable independently of
the chosen reading.

What changes with the reading is solely the \emph{interpretation}
of what the experiment measures. Under R1, an experimental
confirmation of QH1a is an unexplained empirical fact of
coincidence with the sieve structure. Under R2, it is the
empirical expression of an abstract common structure between
sieve and electrodynamics. Under R3, it is a direct measurement
of the ontological informational reality. The three readings are
compatible with the same observations; the choice among them
depends on extra-empirical criteria (parsimony, explanatory power,
metaphysical commitment).

This independence has an important practical consequence for the
experimental evaluation of PT: a test of QH1a can be conducted,
and its results interpreted, without prior commitment on the
ontological status of the $\theta_p$. The theory is testable
independently of the reading. To our knowledge, this is the
first theoretical framework that proposes a structural
dissolution of the Aharonov--Bohm / Bell debate while remaining
empirically falsifiable, and that can be evaluated without first
having to choose between the classical ontological positions.
